% 来源：OEBPS/Text/01_43.xhtml（OPF spine 顺序；仅第一册）
\cleardoublepage
\BookEntry[欧文·斯通《筚路蓝缕的先行者》]{欧文·斯通《筚路蓝缕的先行者》}

\BookCover{images/书单27筚路蓝缕的先行者.jpg}{[美] 欧文·斯通著| Castle Books 2009年版}

\begin{BookSummaryBox}
{\centering \textbf{享誉欧美的传记作家欧文·斯通\\ 记录开发西部的勇士们的故事}\par}
\end{BookSummaryBox}

\subsubsection*{这本书谈什么？}

欧文·斯通是美国作家，他的写作生涯从写剧本开始，后来转向人物传记小说创作，并因此而闻名。他的代表作包括根据梵高的生平撰写的《渴望生活》（1934），根据米开朗琪罗的生平写就的《痛苦与狂喜》（1961）和根据查尔斯·达尔文的生平写成的《起源》（1980）。《筚路蓝缕的先行者》（Men to Match My Mountains）是一部极为优秀的文学作品。这本书描绘了那些19世纪40年代开拓美国西部和加利福尼亚州的勇士的故事。
\par

\subsubsection*{芒格为什么谈起这本书？}

在“论基本的、普世的智慧”的演讲中，芒格谈到基本的文学道理时说：
\par

\begin{quotation}
优秀的文学作品需要读者略加思索才能理解，那样它对读者的影响会很深，你会更牢固地记住它。这就是承诺和保持一贯性的倾向。如果你动脑筋才懂得某个道理，你就会更好地记住它。
\par
\end{quotation}

\subsubsection*{《筚路蓝缕的先行者》作者欧文·斯通如是说}

“生活对于我来说是一次艰难的旅行，我不知道潮水会不会上涨，及至没过嘴唇，甚至涨得更高，但是我要前行。”
\par

“人的童年过得十分缓慢，好像满载的货车；老年的岁月转瞬即逝，犹如夜空的流星。”
\par

“你受到的磨难越多，你就越应该感到高兴。一流的画家就是这样造就出来的。空着肚子比脑满肠肥要强。一颗破碎的心所感受到的不幸，比美满的幸福对你更有好处。永远不要忘记这一点。”
\par
